\chapter{Marked reconstruction, descent and sewing}
\label{ch:vol2-marked-reconstruction}

A boundary junction can detect some boundary objects and miss others.
Its endomorphism algebra determines the objects generated by the
junction. Recovering a larger category requires a generation theorem;
recovering local fields also requires mapping objects that respect
quasi-isomorphisms and geometric composition. These requirements have
different tests.

The simplest test uses the complex
\[
 P_0=\bigl(k[x]\theta\xrightarrow{\ x\ }k[x]\bigr),
 \qquad |\theta|=-1.
\]
It reconstructs the sector supported at \(x=0\), but it cannot detect
the nonzero module \(k[x,x^{-1}]\). For the Ramond junction, the
failure persists without localization: nonzero polynomial boundary
modules with nonnegative energy have zero derived Hom from the junction.
A local residue gives a second obstruction. Its cone belongs to the
stable category generated by the junction but admits no compatible
action of the original differential graded quartet.

We first establish the positive reconstruction statement and retain
its marking. We then prove these obstructions on their full
coefficient categories. The final sections specify the comparisons
needed for collision descent, holomorphic globalization and sewing
with a prepared state. Throughout, complexes are cohomological,
\(M[1]^n=M^{n+1}\), and \(k\) is a field of characteristic zero.

\section{The sector detected by one object}
\label{sec:vol2-generated-sector}

Put \(A=k[x]\) and regard \(P_0\) as a left \(A\)-module complex.
Its augmentation to \(A/(x)\) is a quasi-isomorphism. The derived
category \(D(A)\) is obtained from complexes of left \(A\)-modules
by inverting quasi-isomorphisms. Define its \(x\)-torsion sector by
\begin{equation}\label{eq:vol2-torsion-sector}
 D_{x\text{-tor}}(A)
 =\{M\in D(A):A[x^{-1}]\otimes_A^{\mathbb L}M\simeq0\}.
\end{equation}
This condition is imposed on the derived localization. It does not
impose \(x\)-adic completion.
An object is compact when its derived Hom functor preserves
coproducts. It generates a stable sector when the smallest
full stable subcategory containing it and closed under
coproducts is that sector.

\begin{proposition}\label{prop:vol2-torsion-generator}
The object \(P_0\) is a compact generator of
\(D_{x\text{-tor}}(A)\). It does not generate \(D(A)\).
\end{proposition}
\begin{proof}
Let \(\operatorname{Loc}(P_0)\) be the smallest full stable
subcategory of \(D(A)\) containing \(P_0\) and closed under
coproducts. Localization preserves triangles and coproducts and
kills \(P_0\), so
\(\operatorname{Loc}(P_0)\subset D_{x\text{-tor}}(A)\).
The localization fibre
\[
 \Gamma_xA=[A\longrightarrow A[x^{-1}]]
\]
has its terms in degrees \(0,1\). It is quasi-isomorphic to
\((A[x^{-1}]/A)[-1]\), and
\[
 A[x^{-1}]/A
 =\varinjlim_n(A/x^n,\ \times x).
\]
Each \(A/x^n\) is a finite extension of copies of \(A/x\).
After tensoring with \(M\), each successive factor is a complex
of \(k\)-vector spaces with the quotient \(A\)-action. Such a
complex belongs to \(\operatorname{Loc}(A/x)\): over a field
it is a direct sum of shifts of its cohomology. Consequently
\(\Gamma_xA\otimes_A^{\mathbb L}M\) belongs to
\(\operatorname{Loc}(P_0)\). The localization triangle identifies
this object with \(M\) whenever \(M\) satisfies
\eqref{eq:vol2-torsion-sector}. This proves equality of the two
subcategories. The two-term finite free model makes \(P_0\)
compact.

On the other hand,
\[
 \operatorname{RHom}_A(P_0,A[x^{-1}])
 =[A[x^{-1}]\xrightarrow{\ -x\ }A[x^{-1}]]
\]
is contracted by multiplication by \(-x^{-1}\). Its nonzero
target is therefore invisible to \(P_0\), which excludes
generation of \(D(A)\).
\end{proof}

The general statement uses the same evaluation map. Let \(A\)
now be a dg algebra over a commutative dg ring \(R\), and let
\(P\in D(A)\). Choose a cofibrant representative and write
\[
 E=\operatorname{REnd}_A(P).
\]
We use right \(E\)-modules. Precomposition makes
\(\operatorname{RHom}_A(P,M)\) a right \(E\)-module; the
evaluation action makes \(P\) a left \(E\)-module commuting,
with the graded sign, with its left \(A\)-action. The adjunction is
\begin{equation}\label{eq:vol2-morita-adjunction}
 G(N)=N\otimes_E^{\mathbb L}P,\qquad
 F(M)=\operatorname{RHom}_A(P,M),\qquad G\dashv F.
\end{equation}
Let \(\operatorname{thick}(P)\) denote closure under finite
extensions, shifts and retracts. Let \(\operatorname{Perf}(E)\)
denote the corresponding closure of the free right module \(E\).

\begin{theorem}[Reconstruction of the generated sector]
\label{thm:vol2-marked-morita}
The unit and counit of \eqref{eq:vol2-morita-adjunction} give
inverse equivalences
\[
 \operatorname{thick}(P)\simeq\operatorname{Perf}(E).
\]
If \(P\) is compact in \(D(A)\), they extend to
\[
 \operatorname{Loc}(P)\simeq D(E^{\mathrm{op}}).
\]
Here \(D(E^{\mathrm{op}})\) means the category of right
\(E\)-modules. The image of \(P\) is the free module \(E\);
all specified arrows between objects of the sector are carried
with their full derived mapping spaces.
\end{theorem}
\begin{proof}
There are canonical identifications \(F(P)=E\) and \(G(E)=P\).
The adjunction unit on \(E\) and counit on \(P\) are identities
under these identifications. Both functors are exact and preserve
retracts. The subcategories on which their unit and counit are
equivalences are therefore closed under finite extensions, shifts
and retracts. They contain the respective generators, proving
the first assertion and full faithfulness.

If \(P\) is compact, \(F\) also preserves coproducts; \(G\)
does so as a left adjoint. The same argument now applies to
localizing subcategories. The free right module \(E\) generates
its whole derived category: every module admits a semifree
resolution formed from coproducts of shifts of \(E\), with
successive free attachments. Thus the unit is an equivalence
on every right \(E\)-module. The image of \(G\) is
\(\operatorname{Loc}(P)\), where the counit is an equivalence.
The equivalence is enriched, since the adjunction was formed
from derived tensor and Hom; it therefore carries arrows and
all their homotopies.
\end{proof}

This is the module form of derived Morita theory
\cite{KellerDG2006}. Its domain is part of the theorem. If
\(N\ne0\) and \(F(N)=0\), the counit at \(N\) is \(0\to N\).
Such an \(N\) is a \emph{right orthogonal} to \(P\).
It excludes reconstruction on any category containing \(N\),
without affecting reconstruction on \(\operatorname{thick}(P)\).

\subsection{The unit and the coefficient action}

An equivalence of module categories identifies their centres
when it preserves the specified coefficient action. For ordinary
dg algebras over \(R\), the centre is the Hochschild cochain
\(E_2\)-algebra
\[
 Z_{\mathrm{Mor},R}(A)=\operatorname{RHom}_{A^e}(A,A),
 \qquad A^e=A\otimes_R^{\mathbb L}A^{\mathrm{op}}.
\]
Its \(E_2\) operations arise from composition and the coherent
interchange of endomorphisms of the identity functor.

\begin{proposition}\label{prop:vol2-morita-centre-transport}
An \(R\)-linear derived Morita equivalence identifies these
centres and their actions on modules. If it is represented by
a \(B\)--\(A\) bimodule \(U\) with inverse \(A\)--\(B\)
bimodule \(V\), let \(K\) be an \(R\)-complex. The image of
a bimodule map \(z:K\otimes_R^{\mathbb L}A\to A\) is
\[
 \begin{aligned}
 K\otimes_R^{\mathbb L} B
 &\simeq U\otimes_A^{\mathbb L}
       (K\otimes_R^{\mathbb L} A)\otimes_A^{\mathbb L}V\\
 &\xrightarrow{\,1\otimes z\otimes1\,}
 U\otimes_A^{\mathbb L}V\simeq B.
 \end{aligned}
\]
The comparison uses both inverse bimodule equivalences and
their triangle homotopies.
\end{proposition}
\begin{proof}
A colimit-preserving \(R\)-linear functor on module categories
is determined by its value on the free module. Applying it to
the free bar resolution identifies it with derived tensor by
that value, including the right algebra action and all natural
transformations. In particular, identity-functor endomorphisms
are diagonal-bimodule endomorphisms. Conjugation by the
equivalence and its inverse preserves composition, the
identity functor and the interchange of transformations.
It therefore preserves the centre's \(E_2\) structure.
The displayed formula is this conjugation on evaluation;
tensoring it with a module proves compatibility with its action.
\end{proof}

For example, take \(B=\operatorname{Mat}_2(k)\), \(U=k^2\)
and \(V=(k^2)^*\). The outer product and evaluation give
\[
 U\otimes_k V\xrightarrow{\sim}B,\qquad
 V\otimes_B U\xrightarrow{\sim}k.
\]
The first map sends the four standard tensors to the matrix
units. In the second tensor product, balancing kills the
off-diagonal tensors and identifies the two diagonal tensors.
The inverse images of the units are
\(\sum_i e_i\otimes e_i^*\) and \([e_1^*\otimes e_1]\).
The triangle composites are
\(\sum_i e_i e_i^*(u)=u\) and
\(\sum_i v(e_i)e_i^*=v\).
There is nevertheless no unital algebra map \(B\to k\):
both square-zero off-diagonal units would map to zero,
and so would their products and their sum \(1_B\).
Morita units do not supply scalar augmentations.

\section{Reconstruction with a marked arrow}
\label{sec:vol2-arrow-reconstruction}

An action on an object and an action preserving a chosen
preparation are different data. Fix a closed arrow
\(\nu:Q\to P\) of \(A\)-bimodules, with cofibrant endpoints.
The marking is the entire arrow, including its domain.
Write \(\operatorname{Hom}\) for the bimodule mapping complex
in this subsection. A concrete model for its derived end algebra is
\begin{equation}\label{eq:vol2-arrow-end-algebra}
 E_\nu=\operatorname{End}(Q)\oplus\operatorname{End}(P)
                     \oplus\operatorname{Hom}(Q,P)[-1].
\end{equation}
For a homogeneous triple of degree \(n\), its third component
has degree \(n-1\). Set
\begin{align}
 d(a,b,h)&=(da,db,\nu a-b\nu-dh),
       \label{eq:vol2-arrow-end-d}\\
 (a,b,h)(a',b',h')&=(aa',bb',ha'+(-1)^n bh'),
       \label{eq:vol2-arrow-end-product}\\
 1&=(1_Q,1_P,0).\nonumber
\end{align}
A closed degree-zero triple records two endomorphisms and the
homotopy making the marked square commute.

\begin{proposition}\label{prop:vol2-arrow-end-model}
Equations \eqref{eq:vol2-arrow-end-d} and
\eqref{eq:vol2-arrow-end-product} define a unital dg algebra
model of the endomorphisms of \(\nu\) in the derived arrow
category.
\end{proposition}
\begin{proof}
Put \(F(a,b)=\nu a-b\nu\). The differential is the homotopy-fibre
differential of \(F\), so it squares to zero. The identity
\[
 F(aa',bb')=F(a,b)a'+bF(a',b')
\]
and the Hom differential's Leibniz rule verify the dg Leibniz
identity with the sign in
\eqref{eq:vol2-arrow-end-product}. The third component of
either triple product, for degrees \(n,m\), is
\[
 ha'a''+(-1)^n bh'a''+(-1)^{n+m}bb'h'',
\]
which proves associativity.

To identify the multiplication in the derived category, replace
\(\nu\) by the cofibrant arrow \(i:Q\to T\), where
\[
 \begin{gathered}
 T=P\oplus Q\oplus Q[1],\qquad i(q)=(0,q,0),\\
 d_T(p,q,s)=(d_Pp-\nu s,d_Qq+s,-d_Qs).
 \end{gathered}
\]
The source \(Q\) and the cokernel \(\operatorname{Cone}(-\nu)\)
are cofibrant, which verifies the assertion about the arrow.
The map \(\pi:T\to P\), \(\pi(p,q,s)=p+\nu q\), is a
quasi-isomorphism with contractible kernel and \(\pi i=\nu\).
For a degree-\(n\) triple, define
\[
 \Phi(a,b,h)=(a,B),\qquad
 B(p,q,s)=(bp+(-1)^n hs,\ aq,\ (-1)^n as).
\]
Substitution into the differentials and composition proves that
\(\Phi\) is a strict dg algebra map to \(\operatorname{End}(i)\).
Postcomposition with \((1,\pi)\) identifies its source complex
with \(\operatorname{Hom}(i,\nu)\): the three free coefficients
are \(b,\nu a,(-1)^n h\). Since \(i\) is cofibrant,
\(\operatorname{End}(i)\to\operatorname{Hom}(i,\nu)\)
is a quasi-isomorphism. Thus \(\Phi\) is a quasi-isomorphism,
as required.
\end{proof}

To prescribe the scalar action on the domain of the marking,
form the homotopy pullback of unital \(E_1\)-algebras
\begin{equation}\label{eq:vol2-fixed-mark-classifier}
 Z_\nu=E_\nu\times^h_{\operatorname{REnd}_{A^e}(Q)}R.
\end{equation}
The map from \(R\) is scalar multiplication; the projection
\(Z_\nu\to R\) is an augmentation.

\begin{theorem}\label{thm:vol2-fixed-mark-classifier}
Let \(B\) be an independently specified augmented unital
\(E_1\)-algebra over \(R\), with augmentation \(\varepsilon\).
Then
\[
 \operatorname{Map}_{\operatorname{Alg}_R/R}(B,Z_\nu)
 \simeq \operatorname{Act}^{E_1}_{\varepsilon}(B;\nu).
\]
The right side is the full derived space of actions on the
arrow \(\nu\), together with an identification of the action
on \(Q\) with \(\varepsilon\,1_Q\). It retains the marked
square's homotopy and every higher identification.
\end{theorem}
\begin{proof}
Maps from a cofibrant resolution of \(B\) to the derived end
algebra of an object classify homotopy coherent associative
actions on that object. Applying this construction in the
arrow category gives \(\operatorname{Map}(B,E_\nu)\).
Mapping into the homotopy pullback
\eqref{eq:vol2-fixed-mark-classifier}, and taking the fibre over
the prescribed augmentation, gives
\[
 \operatorname{Map}(B,E_\nu)
 \times^h_{\operatorname{Map}(B,\operatorname{REnd}(Q))}
 \{\varepsilon\,1_Q\}.
\]
This is exactly the space in the statement. The construction
uses derived mapping spaces, so it also identifies their
homotopies in every degree.
\end{proof}

\subsection{Two actions on the same marked object}

The Hochschild centre
\(\mathcal H_A=\operatorname{REnd}_{A^e}(A)\) has two
natural unital \(E_1\) maps to
\(\operatorname{REnd}_{A^e}(P)\):
\begin{equation}\label{eq:vol2-two-tensor-actions}
 \lambda_P(z)=z\otimes_A^{\mathbb L}1_P,\qquad
 \rho_P(z)=1_P\otimes_A^{\mathbb L}z.
\end{equation}
They use the two unit identifications with \(P\).
Derived tensor is an enriched functor, so it preserves
composition and all homotopies. It acts naturally on a supplied
arrow \(\nu\), giving lifts to \(E_\nu\). A lift through
\(Z_\nu\) additionally requires the prescribed scalar action
on \(Q\). No \(E_2\) structure on the end algebra of an arbitrary
object is implied by these \(E_1\) maps.

\begin{example}\label{ex:vol2-split-mark-two-actions}
Let \(A=k[x]\), \(P=A\oplus A^e\), and let \(\nu:A\to P\)
include the diagonal summand. The class \(x\in HH^0(A)\)
acts identically by \(\lambda_P\) and \(\rho_P\) on that
summand. On \(A^e=k[x_L,x_R]\), their difference sends \(1\)
to \(x_L-x_R\ne0\). Since \(P\) is concentrated in degree
zero, this is a nonzero class in
\(H^0\operatorname{REnd}_{A^e}(P)\). Hence the two algebra
maps are not homotopic, even though the marking is split.
\end{example}

The example also explains why the comparison paths in
Chapter~\ref{ch:vol2-currents} cannot be discarded.
An equivalence preserving its polynomial-root object must
preserve the full Massey coset \(-1+pB\). An equivalence over
the level ring must preserve the arrow to the operator complex,
its nonsplit cohomology extension and its derived special
fibre. Retaining only an ordinary scalar centre loses these
data. A fixed-mode gauge equivalence must also retain its
residual automorphisms; equality of orbit sets is weaker.

\section{The Ramond junction and its boundary categories}
\label{sec:vol2-ramond-reconstruction}

We now test generation on a junction defined by its fields.
Let \(R=\mathbb C\llbracket q\rrbracket\) and
\(S_\ell=R/(q^\ell)\), \(\ell\ge1\).
All states over \(S_\ell\) are algebraic.
Completion means \(V\llbracket q\rrbracket\), with a finite
algebraic state at each coefficient; no product over oscillator
energy or internal shifts is taken.

Fix a Ramond coordinate and the normalized coefficient lines.
The boundary algebras have zero differential and regular
self-products:
\[
 A_0=S_\ell[\partial^j\gamma:j\ge0],\qquad
 A_1=\Lambda_{S_\ell}(\partial^j C:j\ge0).
\]
The physical odd field is \(c=qC\). The normalized action
of \(C\) is specified independently, without inverting \(q\).
The junction state complex is
\begin{equation}\label{eq:vol2-ramond-junction}
 P_\ell=S_\ell[x_i,y_n]\otimes\Lambda(\theta_i,\eta_n),
 \qquad
 D=\sum_{i\ge0}x_i\partial_{\theta_i}
             -\sum_{n\ge1}\eta_n\partial_{y_n}.
\end{equation}
The degrees of \(x,y,\theta,\eta\) are \(0,0,-1,1\).
Exterior derivatives are left derivatives. Its two internal
gradings \(W,E\) are
\[
 (W,E)(x_i)=(W,E)(\theta_i)=(1,2i+1),\qquad
 (W,E)(y_n)=(W,E)(\eta_n)=(-1,2n-1).
\]
Here \(E\) is the specified oscillator grading.

The independent boundary actions are
\[
 X_i=x_i,\quad G_n=-\partial_{y_n},\qquad
 K_i=\partial_{\theta_i},\quad H_n=\eta_n.
\]
The even modes commute, and the odd modes supercommute.
Their generating fields are
\begin{align}
 \Gamma(z)&=\sum_{i\ge0}X_i z^{i-1/2}
                  +\sum_{n\ge1}G_nz^{-n-1/2},\label{eq:vol2-ramond-even}\\
 C(z)&=\sum_{n\ge1}H_nz^{n-1/2}
                  +\sum_{i\ge0}K_i z^{-i-1/2}.\label{eq:vol2-ramond-odd}
\end{align}
The even action is originally a right action. Its commutativity
allows left-module notation here; no odd opposite-algebra
identification is made.

Define \(\mathcal C_{0,\ell}\) to consist of internally
bigraded dg modules with this unital even field action such that,
for every vector \(v\), there are \(N,h\) with
\begin{equation}\label{eq:vol2-positive-mode-support}
 (G_n:n>N)v=0,\qquad (G_1,\ldots,G_N)^h v=0.
\end{equation}
Equivalently, \(J_{N,h}v=0\), where
\(J_{N,h}=(G_n:n>N)+(G_1,\ldots,G_N)^h\).
Define \(\mathcal C_{1,\ell}\) using the odd field action and
the condition \(K_iv=0\) for sufficiently large \(i\), for each
\(v\). A finite exterior family is already jointly nilpotent.
All derivative, product and regular substitution identities
are part of these field actions.

The restricted enriched mapping complex is the direct sum
over internal shifts,
\[
 \underline{\operatorname{Hom}}^{\mathrm{res}}(M,N)
 =\bigoplus_{u,v\in\mathbb Z}
   \operatorname{Hom}_{\mathrm{fields,bigr}}
                    (M,N\langle u,v\rangle).
\]
The shift means that a map raises \(W,E\) by \(u,v\).
A degree-\(r\) odd-boundary map obeys
\(f(am)=(-1)^{r|a|}af(m)\); its differential is
\(d_Nf-(-1)^rfd_M\). Derive these mapping complexes at each
\(S_\ell\), then take \(\mathbf R\varprojlim_\ell\).
Weak equivalences are underlying quasi-isomorphisms.

Translation is additional structure: an even operator \(T_M\)
of internal degree \((0,2)\), commuting with \(d_M\), and
\([T_M,Y_M(a,z)]=\partial_zY_M(a,z)\).
Use a superscript \(T\) when it is retained. On \(P_\ell\),
\begin{align}
 T_P={}&\sum_{i\ge0}(i+\tfrac12)
       (x_{i+1}\partial_{x_i}+\theta_{i+1}\partial_{\theta_i})
       \nonumber\\
 &+\sum_{n\ge1}(n+\tfrac12)
       (y_{n+1}\partial_{y_n}+\eta_{n+1}\partial_{\eta_n})
       -\tfrac12y_1x_0+\tfrac12\eta_1\theta_0.
       \label{eq:vol2-ramond-translation}
\end{align}
All sums are finite on an input. The two final terms are
multiplication operators. In particular,
\begin{equation}\label{eq:vol2-preparation-translation}
 T_P1=D(-y_1\theta_0/2)\ne0.
\end{equation}
Both boundary unit fields still act by the identity.
Write \(\mathcal D^T_{a,\ell}\) for the corresponding
category localized at the underlying quasi-isomorphisms.

\subsection{A positive reverse mode comparison}

There is an exact reconstruction of the reverse mode algebra
at the junction. Its recipient is the derived category of
all bigraded modules over
\(\mathsf B_\ell=\Lambda_{S_\ell}(\kappa_i,\eta_n)\),
where \(\kappa_i=K_i\), with the restricted sum of internal
shifts defined above. Write \(P=\varprojlim_\ell P_\ell\)
and \(\mathsf B=\varprojlim_\ell\mathsf B_\ell\) for their
coefficientwise completions.
This paragraph does not impose the additional odd finite-tail
condition on every object.

Put \(a_i=2i+1\), \(b_n=2n-1\). Define an independent
algebra of coefficientwise continuous operators by
\begin{equation}\label{eq:vol2-restricted-reverse-actor}
 (\mathsf A_{\mathrm{res}})^{m,e}
 =
 \left(
 \prod_{|\beta|\le m}
 \ \bigoplus_{\substack{|\alpha|=m-|\beta|\\
                 a\cdot\alpha-b\cdot\beta=e}}
 \mathbb C\,x^\alpha g^\beta
 \right)\llbracket q\rrbracket,\qquad m\ge0.
\end{equation}
Take the coefficientwise direct sum over \(m,e\); all
multiindices have finite support. For fixed \(\beta\),
the sum over \(\alpha\) is finite because the positive
weights \(a_i\) tend to infinity. On an input, only finitely
many \(\beta\) act through \(g_n=-\partial_{y_n}\).
Thus \eqref{eq:vol2-restricted-reverse-actor} is an actual
operator algebra, with composition adding internal degrees.

\begin{proposition}\label{prop:vol2-reverse-mode-equivalence}
The independent right action induces an equivalence
\[
 \mathsf A_{\mathrm{res}}
 \xrightarrow{\sim}
 \underline{\operatorname{REnd}}_{\mathsf B}^{\mathrm{res}}
                                      (P)^{\mathrm{op}}.
\]
It is a statement in this mode category, with the specified
coefficient completion.
\end{proposition}
\begin{proof}
At coefficient level \(S_\ell\), take the semifree resolution
\[
 Q_{\mathsf B}=\mathsf B_\ell[t_i,s_n],\qquad
 d_Q=\sum_i\kappa_i\partial_{t_i}
                         +\sum_n\eta_n\partial_{s_n}.
\]
The new generators have degree zero and internal degrees
\((-1,-a_i)\), \((-1,b_n)\). Use divided monomials.
Its augmentation to \(S_\ell\) is a quasi-isomorphism by
the positive-degree Euler contraction. Filtration by total
\(t,s\)-degree is semifree.

The restricted dual
\(\underline{\operatorname{Hom}}_{\mathsf B}^{\mathrm{res}}
(Q_{\mathsf B},S_\ell)\) is concentrated in degree zero
with zero differential. A map of shift \((m,e)\) can be
nonzero on \(t^{[\alpha]}s^{[\beta]}\) exactly when
\[
 |\alpha|+|\beta|=m,\qquad a\cdot\alpha-b\cdot\beta=e.
\]
Its independent coefficient values give precisely
\eqref{eq:vol2-restricted-reverse-actor} modulo \(q^\ell\).
The primitive coproduct on divided powers identifies
convolution with multiplication of the dual symbols.

To identify the actual action, define the
\(\mathsf B_\ell\)-linear quasi-isomorphism
\[
 f(t^{[\alpha]}s^{[\beta]})
       =\delta_{\alpha,0}(-y)^{[\beta]}
       \quad:Q_{\mathsf B}\longrightarrow P_\ell.
\]
It preserves the constant cohomology. Define a degree-minus-one
map by
\[
 H_i(t^{[\alpha]}s^{[\beta]})
       =\theta_i\delta_{\alpha,0}(-y)^{[\beta]},
 \qquad H_i(bv)=(-1)^{|b|}bH_i(v).
\]
The Clifford identity gives
\[
 d_{\operatorname{Hom}}H_i=x_if-f\partial_{t_i},
 \qquad g_nf=f\partial_{s_n}.
\]
For \(F=x_{i_1}\cdots x_{i_r}g^\beta\), telescope:
\[
 H_F=\sum_{j=1}^r x_{i_1}\cdots x_{i_{j-1}}
 H_{i_j}\partial_{t_{i_{j+1}}}\cdots
             \partial_{t_{i_r}}\partial_s^\beta,
 \qquad dH_F=F_Pf-fF_Q.
\]
The finite-input bounds of
\eqref{eq:vol2-restricted-reverse-actor} make this homotopy
valid for every admitted coefficient sum. Composing with
the augmentation of \(Q_{\mathsf B}\) identifies the actual
action with the corresponding dual coefficient. It induces
an isomorphism on cohomology and was already an associative
derived-algebra map, proving the equivalence.
The coefficient reductions of these Hom models are
surjective, so the ordinary inverse limit computes the
prescribed derived limit. All maps commute with reduction.
\end{proof}

Even this equivalence does not reconstruct all local fields.
The bidegree \((2,0)\) consists of operators
\(\sum_i\lambda_i x_i g_{i+1}\).
Modes of finite differential-polynomial fields in \(\gamma\),
and finite products of such modes, have coefficient sequences
in
\[
 \mathbb C^{(\mathbb N_0)}+\mathbb C[i(i+1)]
\]
at each coefficient of \(q\). Indeed a quadratic derivative
field contributes
\[
 (-1)^{a+b}
 \bigl[(\tfrac12-i)_a(i+\tfrac32)_b
       +(i+\tfrac32)_a(\tfrac12-i)_b\bigr],
\]
where \((z)_a\) is the rising factorial. These polynomials
are invariant under \(i\mapsto-i-1\); the cases \((a,b)=(2r,0)\)
have degree \(2r\) and leading coefficient \(2\), so they
span \(\mathbb C[i(i+1)]\). Products of two linear modes
supply exactly the finitely supported terms. Higher
polynomial degree cannot contribute at weight two.
The sequence \(\lambda_i=2^i\) is absent: every sufficiently
high finite difference kills a polynomial but preserves
\(2^i\). It therefore gives a nonzero action-cone class,
not killed by any power of \(q\). The mode completion and
the local field algebra retain different information.

\section{Polynomial right orthogonals}
\label{sec:vol2-ramond-orthogonals}

Consider the zero-differential single-boundary modules
\[
 N_{0,\ell}=S_\ell[x_i,y_n],\qquad
 N_{1,\ell}=\Lambda_{S_\ell}(\theta_i,\eta_n).
\]
Their modes and gradings are inherited from \(P_\ell\).
Their translations are the even and odd parts, respectively,
of \eqref{eq:vol2-ramond-translation}, including the
multiplication terms \(-y_1x_0/2\) and \(+\eta_1\theta_0/2\).
They satisfy \eqref{eq:vol2-positive-mode-support} or the
odd finite-tail condition. Each has nonnegative energy and
finite free energy components, because all its state generators
have positive energy. Their constant classes are nonzero.

\begin{theorem}\label{thm:vol2-ramond-nongeneration}
In both boundary categories, with translation either retained
or forgotten,
\[
 \underline{\operatorname{RHom}}^{\mathrm{res}}
          (P_\ell,N_{0,\ell})=0,\qquad
 \underline{\operatorname{RHom}}^{\mathrm{res}}
          (P_\ell,N_{1,\ell})=0.
\]
The same holds after the prescribed derived coefficient
completion. Thus the junction does not generate either full
category, even on the objects with nonnegative energy and
finite free energy components.
\end{theorem}

The proof needs the infinite set of modes. A finite cutoff has
a surviving top Koszul class, so its vanishing cannot serve as
a substitute.

\subsection{Cofree reduction}

Put \(B_0=S_\ell[x_i]\) and \(B_1=\Lambda_{S_\ell}(\eta_n)\).
Let \(C_+\) be the coalgebra with finite-support divided-power
basis \(e_\alpha\) and
\[
 \Delta e_\alpha=\sum_{\beta+\gamma=\alpha}
                         e_\beta\otimes e_\gamma.
\]
Let \(C_\theta=\Lambda_{S_\ell}(\theta_i)\) be the exterior
coalgebra with primitive generators. There are exact
presentations
\begin{equation}\label{eq:vol2-supported-comodule-categories}
 \begin{aligned}
 \mathcal C_{0,\ell}&\simeq
 C_+\text{-}\operatorname{Comod}(B_0\text{-}\operatorname{Mod}_{\mathrm{bigr}}),\\
 \mathcal C_{1,\ell}&\simeq
 C_\theta\text{-}\operatorname{Comod}(B_1\text{-}\operatorname{Mod}_{\mathrm{bigr}}).
 \end{aligned}
\end{equation}
Indeed, the even coaction
\(\rho(v)=\sum_\alpha e_\alpha\otimes G^\alpha v\) is finite
precisely by \eqref{eq:vol2-positive-mode-support}. The odd
coaction uses finite subsets and the exterior signs.
Conversely these modes define the fields of every differential
polynomial. Move annihilators to the right: their indices and
powers are bounded on an input, and a fixed Laurent coefficient
then bounds the sum of the remaining creation indices.
Every coefficient is finite. Zero self-contractions give the
product and substitution identities. Intertwining all modes
is equivalent to intertwining these fields, proving
\eqref{eq:vol2-supported-comodule-categories}.

The forgetful functor \(U\) is exact. Its cofree right adjoint
\(F(V)=C\otimes_{S_\ell}V\) is exact because either coalgebra
is a direct sum of shifts of free coefficient modules.
For a K-injective base resolution \(I\) of \(V\) and an
acyclic comodule \(Z\),
\[
 \operatorname{Hom}_{\mathcal C}(Z,FI)
 =\operatorname{Hom}_{B}(UZ,I)
\]
is acyclic. Hence \(FI\) computes the derived cofree object,
including for unbounded complexes. In every internal shift,
\begin{equation}\label{eq:vol2-cofree-derived-adjunction}
 \operatorname{RHom}_{\mathcal C}(M,FV)
 \simeq\operatorname{RHom}_{B}(UM,V).
\end{equation}
Taking the exact direct sum gives the restricted enrichment.

The targets are \(F(B_0)\) and \(F(B_1)\); on the even side
use \(e_\alpha=(-y)^\alpha/\alpha!\).
The appropriate Koszul or polynomial de Rham factor of
\(UP_\ell\) contracts by the Euler homotopy, leaving \(S_\ell\).
For instance, over \(B_0\) set \(y,\eta\) to zero and retain
\(B_0\otimes\Lambda(\theta_i)\), \(d\theta_i=x_i\);
over \(B_1\) set \(x,\theta\) to zero and retain
\(B_1[y_n]\), \(dy_n=-\eta_n\).
Thus it remains to prove
\begin{equation}\label{eq:vol2-infinite-derived-duals}
 \underline{\operatorname{RHom}}_{B_0}^{\mathrm{res}}(S_\ell,B_0)=0,
 \qquad
 \underline{\operatorname{RHom}}_{B_1}^{\mathrm{res}}(S_\ell,B_1)=0.
\end{equation}

\subsection{A finite head in every cohomological degree}

For \(B_0\), use the semifree resolution
\(B_0\otimes\Lambda(\theta_i)\) with \(d\theta_i=x_i\).
Its positive-degree Euler contraction proves exactness.
Semifreeness gives K-projectivity: against an acyclic target,
a primitive of a closed Hom map is chosen successively on
free generators in each internal shift.

A degree-\(p\) Hom cochain is an alternating family of finite
polynomials indexed by \(p\)-element subsets; the family itself
need not have finite support. With dual exterior symbols
\(\epsilon_i\), its raw differential is
\((-1)^{p+1}\sum_i x_i\epsilon_i\wedge-\).
Conjugation in degree \(p\) by \((-1)^{p(p+1)/2}\) changes it to
\[
 \delta=\sum_i x_i\epsilon_i\wedge-.
\]
Each output coefficient uses only the finitely many indices
in its indexing subset. Define
\[
 Q_jf=\frac{f-f|_{x_j=0}}{x_j},\qquad
 h_j=\iota_{\epsilon_j}Q_j.
\]
Let \(\pi_j\) retain terms containing \(\epsilon_j\) and
set \(x_j=0\) in their coefficients. Polynomial division gives
\begin{equation}\label{eq:vol2-head-contraction}
 \delta h_j+h_j\delta=1-\pi_j.
\end{equation}
On a term without \(\epsilon_j\), use \(Q_j(x_jf)=f\).
On a term with it, the two sides subtract the constant
coefficient in \(x_j\). Cross-index terms anticommute.

For a head of length \(m\), put
\[
 H_m=h_0+\pi_0h_1+\cdots+\pi_0\cdots\pi_{m-2}h_{m-1},
 \qquad \Pi_m=\pi_0\cdots\pi_{m-1}.
\]
The projections commute with \(\delta\), so telescoping gives
\begin{equation}\label{eq:vol2-head-telescope}
 \delta H_m+H_m\delta=1-\Pi_m.
\end{equation}
On a degree-\(p\) cocycle choose \(m=p+1\).
The image of \(\Pi_m\) requires \(m\) exterior symbols and
therefore vanishes. This constructs a primitive of every
cocycle. Each primitive still has finite polynomial output
on each input and preserves the two internal shifts.

For \(B_1\), use
\[
 Q=B_1\otimes S_\ell[s_1,s_2,\ldots],\qquad
 d_Q=\sum_n\eta_n\partial_{s_n},\qquad |s_n|=0,
\]
with divided monomials and
\((W,E)(s_n)=(-1,2n-1)\). The positive-degree Euler contraction
and filtration by total \(s\)-degree prove that this is a
semifree resolution. A degree-\(p\) Hom cochain is a family
\(F_\alpha\in B_1^p\). Its raw differential is
\(-\sum_{\alpha_n>0}\eta_nF_{\alpha-e_n}\):
the Hom sign and graded left-linearity cancel their
\((-1)^p\) factors. Conjugating by \((-1)^p\) gives
\[
 (\delta F)_\alpha
       =\sum_{\alpha_n>0}\eta_nF_{\alpha-e_n}.
\]
Now set
\[
 (h_jF)_\alpha=\iota_{\eta_j}F_{\alpha+e_j},\qquad
 (\pi_jF)_\alpha=
 \begin{cases}
 \eta_j\wedge\iota_{\eta_j}F_\alpha,&\alpha_j=0,\\
 0,&\alpha_j>0.
 \end{cases}
\]
The exterior contraction identity proves
\eqref{eq:vol2-head-contraction}, separately for
\(\alpha_j=0\) and \(\alpha_j>0\).
The same telescoping proof applies to the head \(1,\ldots,m\).
Its top projection forces \(m\) distinct exterior output
generators, hence vanishes for \(m=p+1\). These homotopies
query finitely many input coefficients and preserve both
internal shifts. This proves
\eqref{eq:vol2-infinite-derived-duals}.

\subsection{Translation and completion}

Let \(\delta_T\) be the derivation of boundary modes induced
by translation. Adjoin \(T\) by \(Ta-aT=\delta_T(a)\).
Induction from the category forgetting translation has
underlying module \(S_\ell[T]\otimes M\), and
\[
 a(T^r\otimes v)=
 \sum_{j=0}^r(-1)^j\binom rj
 T^{r-j}\otimes\delta_T^j(a)v.
\]
It remains inside the admitted support class. If
\(J_{N,h}\) kills \(v\), then \(J_{N+r,h+r}\) kills
\(T^r\otimes v\): at most \(r\) annihilators can be
differentiated, leaving at least \(h\) original ones.
Their indices decrease by at most \(r\), and the even creation
modes at the seam commute with them. The odd tail bound
follows from the same finite index change.

For a translated module, the relative sequence
\[
 \begin{gathered}
 0\longrightarrow\operatorname{Ind}(UM)^{(E+2)}
 \xrightarrow{\ \partial\ }\operatorname{Ind}(UM)
 \longrightarrow M\longrightarrow0,\\
 \partial(a\otimes m)=aT\otimes m-a\otimes T_Mm
 \end{gathered}
\]
is exact by division by the monic operator \(T-T_M\).
The first term's generators have energy increased by two.
The covariance identity makes the map balanced; commutation
with the differential makes it a chain map.
Derived adjunction therefore computes translated Hom by a
two-term totalization of underlying derived Hom objects.
Both terms vanish by the preceding calculation. Finally
\(\mathbf R\varprojlim_\ell\) of the zero mapping objects
is zero, while \(N_a\) retains its constant class.
This completes the proof of
Theorem~\ref{thm:vol2-ramond-nongeneration}.

The proof excludes generation in the stated categories.
It does not calculate a dual for an unspecified native tensor.
Regular ambient mode modules and supported boundary modules
also have different admissibility conditions. An assertion
about one of their derived duals cannot be transferred to the
other without a comparison.

\section{A local residue whose cone has no quartet lift}
\label{sec:vol2-quartet-cone}

The polynomial right orthogonals have zero differential.
They cannot carry the original dg quartet, since its relations
would require nonzero boundary actions to be nullhomotopic.
Imposing quartet liftability on each object might therefore
appear to select a reconstructible sector. The next calculation
shows that this condition is not stable under the boundary
cones already generated by \(P_\ell\).

On \(P_\ell\), the original conjugate fields are
\[
 \begin{aligned}
 \widehat\beta(z)&=\sum_{n\ge1}y_nz^{n-1/2}
                     +\sum_{i\ge0}\partial_{x_i}z^{-i-1/2},\\
 b(z)&=\sum_{i\ge0}\theta_i z^{i-1/2}
                     +\sum_{n\ge1}\partial_{\eta_n}z^{-n-1/2}.
 \end{aligned}
\]
A compatible quartet extension on another boundary object means
fields with the same normalized Weyl and Clifford relations,
all original zero mixed supercommutators, and
\[
 [d,b(z)]=\Gamma(z),\qquad [d,\widehat\beta(z)]=-C(z).
\]
In particular its mode \(b_0\) is even-boundary-linear and
satisfies \([d,b_0]=X_0\). Thus the \(X_0\)-action vanishes
in the derived even-boundary category after forgetting
translation. This necessary condition is invariant under
boundary quasi-isomorphism.

Define the degree-two operator
\begin{equation}\label{eq:vol2-ramond-residue}
 \Omega=\sum_{i\ge0}(2i+1)\eta_{i+1}\partial_{\theta_i}.
\end{equation}
It is finite on each input, has internal degree \((-2,0)\),
and commutes with \(D\) and the even boundary modes.

\begin{lemma}\label{lem:vol2-residue-translation}
The operator \(\Omega\) is a translation-invariant local
residue:
\[
 \Omega=-\operatorname{Res}_{z=0}C(z)\partial_z C(z),
 \qquad [T_P,\Omega]=0.
\]
Within the family \(\sum_i a_i\eta_{i+1}\partial_{\theta_i}\),
translation invariance is equivalent to
\(a_i=(2i+1)a_0\).
\end{lemma}
\begin{proof}
A residue contains one creation mode and one annihilation
mode. The pair with index \(i\) contributes
\[
 (i+\tfrac12)\partial_{\theta_i}\eta_{i+1}
 -(i+\tfrac12)\eta_{i+1}\partial_{\theta_i}
 =-(2i+1)\eta_{i+1}\partial_{\theta_i}.
\]
The field \(C\) has zero self-contraction, so this is a mode
of its differential-polynomial local field.
Translation gives
\[
 \begin{gathered}
 [T_P,\partial_{\theta_0}]=\eta_1/2,\qquad
 [T_P,\partial_{\theta_i}]
       =(\tfrac12-i)\partial_{\theta_{i-1}}\quad(i\ge1),\\
 [T_P,\eta_n]=(n+\tfrac12)\eta_{n+1}.
 \end{gathered}
\]
The coefficient of
\(\eta_{i+2}\partial_{\theta_i}\) in the commutator is
\(a_i(i+\tfrac32)-a_{i+1}(i+\tfrac12)\).
These operators are independent, as evaluation on \(\theta_i\)
shows. Their vanishing is the stated recurrence. The remaining
seam is \(a_0\eta_1^2/2=0\).
\end{proof}

Let \(a\in S_\ell\), and adjoin one exterior variable
\(\epsilon\) with degree \(-1\) and internal degree \((2,0)\).
Set
\begin{equation}\label{eq:vol2-residue-cone}
 \mathcal E_{\ell,a}=P_\ell\otimes\Lambda(\epsilon),
 \qquad d_a=D+a\epsilon\Omega,\qquad
 T_{\mathcal E}=T_P\otimes1.
\end{equation}
The even boundary modes are unchanged. Equivalently,
\[
 d_a\theta_i=x_i+a(2i+1)\epsilon\eta_{i+1},\qquad
 d_a y_n=-\eta_n,\qquad
 d_a x_i=d_a\eta_n=d_a\epsilon=0.
\]
The differential squares to zero, commutes with translation,
and preserves the support and internal gradings. The short
exact sequence
\[
 0\longrightarrow\epsilon P_\ell
 \longrightarrow\mathcal E_{\ell,a}
 \longrightarrow P_\ell\longrightarrow0
\]
has connecting morphism \(a\Omega\). Its first term is a
cohomological shift of \(P_\ell\), with weight increased by two.
Thus \(\mathcal E_{\ell,a}\) belongs to the stable closure of
the admitted shifts of \(P_\ell\).
The new variable has energy zero and is exterior, so each
energy component is the direct sum of two finite free
components with different internal shifts. Energy remains
nonnegative.

Its underlying coefficient cohomology is
\(\Lambda_{S_\ell}(\epsilon)\). Indeed, the homogeneous
polynomial change
\[
 x_i'=x_i+a(2i+1)\epsilon\eta_{i+1}
\]
leaves independent pairs \(d\theta_i=x_i'\) and
\(dy_n=-\eta_n\). The Euler homotopy contracts all positive
pair-degree, with \(\epsilon\) a spectator.
This change is not linear for the original boundary action.
To test that action, one must retain a boundary coordinate.

\begin{proposition}\label{prop:vol2-residue-derived-detector}
Forget translation and restrict to the two even modes
\(x=X_0\) and \(g=G_1\). The derived functor
\[
 \mathcal L_g(M)
   =M\otimes^{\mathbb L}_{S_\ell[g]}S_\ell,\qquad g\mapsto0,
\]
retains the \(x\)-action and satisfies
\[
 H^*(\mathcal L_g\mathcal E_{\ell,a})
       \cong e\,\Lambda_{S_\ell}(\epsilon,\eta_1),
 \qquad |e|=-1.
\]
In the ordered basis
\((e,\epsilon e,\eta_1e,\epsilon\eta_1e)\), \(x\) has matrix
\begin{equation}\label{eq:vol2-residue-coordinate-matrix}
 X_a=
 \begin{pmatrix}
 0&0&0&0\\
 0&0&0&0\\
 0&0&0&0\\
 -a&0&0&0
 \end{pmatrix}.
\end{equation}
It is nonzero whenever \(a\ne0\).
\end{proposition}
\begin{proof}
The bounded free resolution
\([S_\ell[g]e\xrightarrow{g}S_\ell[g]]\), with
\((W,E)(e)=(1,-1)\), computes the functor even on
unbounded complexes. On \(\mathcal E_{\ell,a}\),
\(g=-\partial_{y_1}\) is degreewise surjective and commutes
with \(d_a\). Its two-term Koszul tensor is therefore
quasi-isomorphic to its kernel in the indicated shift.
This follows from the short exact sequence of complexes;
polynomial integration need not be a chain map.

The kernel contains precisely the states independent of
\(y_1\). The coordinates \(x_i'\) contract every
\((\theta_i,x_i')\) pair, and the remaining
\((y_n,-\eta_n)\) pairs for \(n\ge2\). What survives is
\(\Lambda(\epsilon,\eta_1)\). Since
\(x_0=x_0'-a\epsilon\eta_1\), its action is exactly
\eqref{eq:vol2-residue-coordinate-matrix}.

For an explicit last contraction, put
\[
 B=\Lambda_{S_\ell}(\epsilon,\eta),\qquad
 K_a=B[x]\otimes\Lambda(\theta),\qquad
 d\theta=x+a\epsilon\eta.
\]
The multiplier \(x+a\epsilon\eta\) is monic, hence injective
on \(B[x]\), including over the nonreduced ring \(S_\ell\).
Polynomial division gives
\[
 Q_a(u)=\frac{u-u|_{x=-a\epsilon\eta}}{x+a\epsilon\eta},
 \qquad h(u)=\theta Q_a(u),\qquad h(\theta u)=0.
\]
With \(\rho\) setting \(\theta=0,x=-a\epsilon\eta\) and
\(\iota\) including exterior representatives,
\(dh+hd=1-\iota\rho\). Thus the four surviving classes are
free over \(S_\ell\), and a nonzero entry \(-a\) cannot
vanish by coefficient torsion.
\end{proof}

\begin{theorem}\label{thm:vol2-no-stable-quartet-sector}
No full replete stable subcategory of
\(\mathcal D^T_{0,\ell}\), containing \(P_\ell\) and
closed under its admitted internal shifts, can consist
entirely of objects admitting the original dg quartet
extension, even up to boundary quasi-isomorphism.
The assertion survives coefficient completion and any further
support restrictions on the extra quartet operators.
\end{theorem}
\begin{proof}
Fullness admits the translated local morphism \(\Omega\)
between shifts of \(P_\ell\); stability admits
\(\mathcal E_{\ell,1}\). Any quartet lift would make its
\(x_0\)-action zero by the boundary-linear homotopy \(b_0\).
Restriction to two modes and derived tensor send a zero
morphism to zero. Proposition~\ref{prop:vol2-residue-derived-detector}
contradicts this. No support condition on \(b_0\) was used.
The coefficient \(-1\) survives already modulo \(q\), so
reduction proves the completed assertion.
\end{proof}

More precisely, \(\mathcal E_{\ell,a}\) is quartet-liftable
if and only if \(a=0\). At \(a=0\) the original operators act
on the first tensor factor. Every nonzero \(a\) is excluded
by the same matrix. The underlying vector complex still has
the primitive
\[
 d_a(\theta_0-a\epsilon y_1)=x_0,
 \qquad
 [g_1,\theta_0-a\epsilon y_1]=a\epsilon.
\]
The second equation is the failure of boundary-linearity.

The integral physical residue is
\[
 -\operatorname{Res}c(z)\partial_zc(z)=q^2\Omega.
\]
Its cone is obstructed for \(\ell\ge3\); for \(\ell=1,2\)
it splits. As a map of complex vector spaces,
\(\operatorname{rank}X_{q^2}=\max(\ell-2,0)\).
Thus the complete lattice is obstructed already modulo
\(q^3\), without cancelling or inverting \(q\).

\subsection{The extension class and its local contact}

The two-mode class has a direct Ext interpretation. Over
\(A=S_\ell[x,g]\), let
\[
 P^{(2)}=S_\ell[x,y]\otimes\Lambda(\theta,\eta),
 \qquad D\theta=x,\quad Dy=-\eta,\quad g=-\partial_y.
\]
The finite free resolution of its constant cohomology is
\[
 Q=A\otimes\Lambda(u,v),\qquad
 |u|=|v|=-1,\quad du=x,\quad dv=g.
\]
The \(A\)-linear map \(f:Q\to P^{(2)}\) given by
\(f(1)=1,f(u)=\theta,f(v)=f(uv)=0\) is a quasi-isomorphism.
For \(\omega=\eta\partial_\theta\), define the degree-one
map \(H\) by \(H(u)=-y\) and zero on the other generators.
Then
\[
 (d_{\operatorname{Hom}}H)(u)=\eta,\qquad
 (d_{\operatorname{Hom}}H)(uv)=-1.
\]
Consequently \(\omega f-d_{\operatorname{Hom}}H\) sends
\(uv\) to \(1\) and the other basis elements to zero.
It is the image under \(f\) of the degree-two endomorphism
of \(Q\) with that same formula. Compose this endomorphism
with the augmentation \(Q\to S_\ell\). The resulting
Hom complex has zero differential, and the value \(1\)
on \(uv\) is its top generator. Thus the class generates
\(\operatorname{Ext}^2_A(S_\ell,S_\ell)=S_\ell\).
The discarded all-mode residue tail has the explicit
two-mode-linear homotopy
\[
 \left[D,-\sum_{i\ge1}(2i+1)y_{i+1}\partial_{\theta_i}\right]
   =\sum_{i\ge1}(2i+1)\eta_{i+1}\partial_{\theta_i}.
\]
Every sum is finite on an input. Thus the Ext class is the
restriction of the actual all-mode residue.

There is a local correction of the field differential, but
it changes a mixed contact. Put
\[
 b_a(z)=b(z)-2a\epsilon\partial_z\widehat\beta(z).
\]
Oscillator commutators give
\[
 [\Omega,b(z)]=2\partial_zC(z),\qquad
 [\Omega,\widehat\beta(z)]=0.
\]
The tensor sign gives
\([d_a,\epsilon\partial_z\widehat\beta]=\epsilon\partial_zC\).
Hence
\[
 [d_a,b_a]=\Gamma,\qquad
 [d_a,\widehat\beta]=-C,\qquad [T_{\mathcal E},b_a]=\partial_zb_a.
\]
However, if
\(\Delta(z,u)=\sum_{r\in\mathbb Z}z^{r-1/2}u^{-r-1/2}\),
then
\begin{equation}\label{eq:vol2-quartet-derivative-contact}
 [\Gamma(u),b_a(z)]=2a\epsilon\,\partial_z\Delta(z,u).
\end{equation}
For example its modes are
\[
 [g_n,b_{a,-i}]=a(2i+1)\epsilon\delta_{n,i+1},\qquad
 [x_i,b_{a,n}]=-a(2n-1)\epsilon\delta_{i,n-1}.
\]
This is a local derivative contact. It is not the original
zero mixed bracket, and
Theorem~\ref{thm:vol2-no-stable-quartet-sector} excludes
every correction restoring all original relations on this
derived boundary object.

\section{Local fields must respect derived equivalence}
\label{sec:vol2-polynomial-fields-derived}

The mapping object is as consequential as the object class.
Local polynomial endomorphism fields at a fixed junction do
not automatically extend to a derived bifunctor on all pairs.
The obstruction is visible on the acyclic localization
\[
 Q_\ell=P_\ell[x_0^{-1}].
\]
The variables retain their gradings, with
\((W,E)(x_0^{-1})=(-1,-1)\).
Laurent differentiation extends all quartet operators and
translation. The boundary support conditions still hold:
inverting \(x_0\) changes neither finite \(y\)-degree nor
finite \(\theta\)-support. The multiplication operator
\(h=m_{\theta_0/x_0}\) satisfies
\[
 Dh+hD=1
\]
as an operator anticommutator. Thus \(Q_\ell\to0\) is a
quasi-isomorphism preserving all specified structures.
The contracting homotopy need not preserve them.

Fix the following polynomial field prescription. A normal
symbol uses the original creation and annihilation variables,
with polynomial creation coefficients and a uniform bound on
total symbol degree at each coefficient of \(q\). It can
contain infinitely many mode monomials, provided its value
on each input is a field with finite lower Laurent order and
honest state coefficients. Inverse creation variables are
not among its generators. Impose quartet locality,
translation covariance and exact linearity for the selected
boundary. For an endomorphism, locality means the first
equation below for each quartet field \(f\) and some finite
order \(K\); covariance means the second:
\[
 (u-z)^K[f(u),F(z)]=0,\qquad [T,F(z)]=\partial_zF(z),
\]
together with supercommutation with every boundary mode.
For \(F:P_\ell\to Q_\ell((z^{1/2}))\), use the corresponding
rectangular equations and differential
\(DF-(-1)^{|F|}FD\).
Call these complexes \(\mathcal E_{a,\mathrm{poly}}^T(M)\)
and \(\mathcal H_{a,\mathrm{poly}}^T(P_\ell,Q_\ell)\).

\begin{proposition}\label{prop:vol2-localization-polynomial-fields}
Extension and restriction identify
\[
 \mathcal E_{a,\mathrm{poly}}^T(Q_\ell)
 \cong\mathcal E_{a,\mathrm{poly}}^T(P_\ell)
 \cong\mathcal H_{a,\mathrm{poly}}^T(P_\ell,Q_\ell)
 \qquad(a=0,1).
\]
The identity of \(Q_\ell\) and inclusion
\(\iota:P_\ell\to Q_\ell\) represent nonzero degree-zero
classes in their respective complexes.
\end{proposition}
\begin{proof}
A normal expression has bounded order in
\(\partial_{x_0}\). On a fixed Laurent input, only finitely
many other annihilator rows can act, by the input's finite
support. Each such row has a lower Laurent bound in the field
variable and finite polynomial coefficients. The expression
therefore extends to \(Q_\ell\).

Successive evaluation on polynomial monomials recovers each
normal row: subtract the proper subrows already recovered,
then divide by the monomial factorial. This proves
faithfulness over \(S_\ell\). Every polynomial expression on
\(Q_\ell\) preserves \(P_\ell\); every rectangular polynomial
expression lands in \(P_\ell\). The field equations are
identities of these bounded-degree polynomial operators,
coefficient by coefficient, so they extend and restrict
faithfully.

Let \(\varepsilon:P_\ell\to S_\ell\) extract the scalar
coefficient. Then \(D1=0\) and \(\varepsilon D=0\).
If a degree-minus-one polynomial field had \(dF=\iota\),
its value on \(1\) would lie in \(P_\ell((z^{1/2}))\).
Applying \(\varepsilon\) to \(D(F(z)1)=1\) would give
\(0=1\). The diagonal identity is excluded in the same
way by restriction.
\end{proof}

\begin{theorem}\label{thm:vol2-polynomial-fields-not-derived}
No reduced bifunctor invariant under underlying
quasi-isomorphisms can be quasi-isomorphic to the unchanged
polynomial field complex on \((P_\ell,Q_\ell)\), or on
every diagonal.
Here reduced means that its value is zero if either input
is zero. The obstruction holds with translation and after
coefficient completion.
\end{theorem}
\begin{proof}
Derived invariance and \(Q_\ell\simeq0\) force both
cross-object and diagonal values involving \(Q_\ell\) to
vanish. Proposition~\ref{prop:vol2-localization-polynomial-fields}
contradicts the proposed comparisons. A completed primitive
would reduce to an impossible primitive modulo \(q\).
\end{proof}

This does not exclude a derived replacement agreeing with
the polynomial actor only at \(P\). It also does not apply
to a recipient imposing nonnegative energy on every object:
\(Q_\ell\) contains \(x_0^{-m}\), and its energy-zero part
contains the independent states
\(x_i x_0^{-(2i+1)}\). The polynomial right orthogonals of
Theorem~\ref{thm:vol2-ramond-nongeneration} require neither
of these features.

\subsection{An inverse field and a remaining odd obstruction}

An enlarged local field class supplies an even-boundary
primitive explicitly. Write
\[
 a(z)=\sum_{i\ge0}x_i z^i,\qquad
 p(z)=\sum_{n\ge1}\partial_{y_n}z^{-n},\qquad
 \Gamma(z)=z^{-1/2}(a(z)-p(z)).
\]
Since \(x_0\) is invertible, set
\begin{equation}\label{eq:vol2-local-inverse-field}
 I(z)=z^{1/2}\sum_{r\ge0}a(z)^{-r-1}p(z)^r,\qquad
 H(z)=b(z)I(z).
\end{equation}

\begin{proposition}\label{prop:vol2-local-inverse-homotopy}
These are fields with finite Laurent-polynomial state
coefficients on every input. They satisfy
\[
 \Gamma I=I\Gamma=1,\qquad [D,H]=1.
\]
The field \(H\) is translation-covariant, local with the
quartet and exactly even-boundary-linear.
There is no degree-minus-one strictly odd-boundary-linear
finite-output field \(P_\ell\to Q_\ell((z^{1/2}))\) with
differential \(\iota\), nor such an endomorphism field of
\(Q_\ell\) with differential \(1\).
\end{proposition}
\begin{proof}
On a vector of \(y\)-degree \(d\), only \(r\le d\) contributes.
At a fixed coefficient of \(z\),
\[
 a(z)^{-s}=x_0^{-s}
       \left(1+\sum_{i\ge1}(x_i/x_0)z^i\right)^{-s}
\]
has finitely many partition terms. The geometric-series
telescoping is therefore finite on every input and proves
the inverse identities. The factors \(b,I\) have zero
mutual contractions; applying inverse-field coefficients
does not increase \(y\)-degree or \(\eta\)-support.
Their product is a field, with \([D,I]=0\), and
\([D,H]=[D,b]I=\Gamma I=1\).

Put \(J=(a-p)^{-1}\). Modewise differentiation gives
\[
 [\partial_{x_i},J]=-z^iJ^2,\qquad
 [y_n,J]=-z^{-n}J^2.
\]
With
\(\Delta(u,z)=\sum_{r\in\mathbb Z}u^{r-1/2}z^{-r-1/2}\),
the nonzero commutators are
\[
 \begin{aligned}
 [\widehat\beta(u),H(z)]&=-\Delta(u,z)b(z)I(z)^2,\\
 [C(u),H(z)]&=\Delta(u,z)I(z).
 \end{aligned}
\]
Each mode coefficient is finite on an input, and
\((u-z)\Delta(u,z)=0\) proves locality.
Exact commutation with \(\Gamma\) proves even-boundary
linearity. Differentiating \(\Gamma I=1\), and commuting it
with \(T\), gives
\([T,I]=-I(\partial_z\Gamma)I=\partial_zI\), then the
covariance of \(H\).

For the odd obstruction, every \(K_i\) kills \(1\).
Strict graded linearity of a degree-minus-one field gives
\(K_iF(z)1=0\) for every \(i\). The common kernel of the
\(K_i=\partial_{\theta_i}\) is
\[
 S_\ell[x_i,x_0^{-1},y_n]\otimes\Lambda(\eta_n),
\]
which has no negative cohomological degree. Thus \(F(z)1=0\),
and its differential cannot send \(1\) to \(1\).
No locality or translation assumption is needed for this
last argument.
\end{proof}

The enlargement in \eqref{eq:vol2-local-inverse-field} is
substantial. On the constant state,
\[
 \begin{gathered}
 H(z)1=\frac{\sum_i\theta_i z^i}{\sum_i x_i z^i},\\
 H_0(1)=\theta_0/x_0,\quad
 H_1(1)=\theta_1/x_0-\theta_0x_1/x_0^2.
 \end{gathered}
\]
Already the \(\theta_0,x_0,x_1\) terms have pole order
\(m+1\) at coefficient \(z^m\). They are honest finite
states coefficientwise, with no uniform pole or symbol-degree
bound. Enlarging symbols can repair the even primitive;
it cannot repair strict odd-boundary linearity.

\section{What full quartet linearity does reconstruct}
\label{sec:vol2-full-quartet-category}

There is a positive reconstruction with a different mapping
category. Let \(\mathscr W_\ell\) be the algebra of finite
polynomial oscillator operators
\[
 x_i,y_n,\theta_i,\eta_n,\quad
 \xi_i=\partial_{x_i},\quad\zeta_n=\partial_{y_n},
 \quad\kappa_i=\partial_{\theta_i},\quad\pi_n=\partial_{\eta_n},
\]
with the normalized Weyl and Clifford relations. All other
supercommutators vanish. Its differential \(\delta=[D,-]\)
satisfies
\[
 \delta\theta_i=x_i,\qquad \delta y_n=-\eta_n,\qquad
 \delta\xi_i=-\kappa_i,\qquad\delta\pi_n=-\zeta_n,
\]
and kills the other generators.

\begin{proposition}\label{prop:vol2-full-quartet-morita}
The unit inclusion \(S_\ell\to\mathscr W_\ell\) and
evaluation \(\mathscr W_\ell\to P_\ell\), \(w\mapsto w1\),
are quasi-isomorphisms. For full quartet-linear morphisms,
\[
 D(\mathscr W_\ell)\simeq D(S_\ell),\qquad
 \operatorname{REnd}_{\mathscr W_\ell}(P_\ell)\simeq S_\ell.
\]
In the bigraded category, the internal shifts of \(P_\ell\)
form the compact generating family.
\end{proposition}
\begin{proof}
Normal ordering gives a free superpolynomial symbol basis.
The differential preserves this order and has exactly the
four displayed pair differentials. In positive total symbol
degree \(N\), the homotopy is
\[
 h=N^{-1}\left(
 \sum_i\theta_i\partial_{x_i}-\sum_n y_n\partial_{\eta_n}
 -\sum_i\xi_i\partial_{\kappa_i}-\sum_n\pi_n\partial_{\zeta_n}
 \right).
\]
Its anticommutator with \(\delta\) is the identity in positive
degree. It is finite on each symbol and preserves internal
degrees. Only the identity class survives; evaluation on \(1\)
identifies it with the constant cohomology of \(P_\ell\).

Tensoring the coefficient contraction with any complex \(V\)
shows that \(V\to\mathscr W_\ell\otimes_{S_\ell}V\) is a
quasi-isomorphism. For a dg \(\mathscr W_\ell\)-module \(M\),
the multiplication map
\(\mathscr W_\ell\otimes_{S_\ell}M\to M\) is a module map
whose composite with the coefficient unit is the identity.
It too is a quasi-isomorphism. These are the derived unit
and counit for extension and restriction along
\(S_\ell\to\mathscr W_\ell\).
The assertions about the end and compact generators follow.
\end{proof}

The scalar symbol projection used in the proof is not an
algebra augmentation. Indeed,
\(\xi_i x_i=x_i\xi_i+1\) has scalar part \(1\), while
\(\xi_i\) and \(x_i\) each have scalar part zero.
The multiplicative quasi-isomorphism is the unit inclusion.

One may also impose polynomial joint support: on every
vector all but finitely many annihilators vanish, and the
remaining even annihilators are jointly locally nilpotent.
For one even pair \([a,z]=1\), the finite-on-input operator
\[
 \Pi_{z,a}=\sum_{r\ge0}\frac{(-1)^r}{r!}z^ra^r
\]
projects to \(\ker a\); for an odd pair \(az+za=1\),
the projector is \(1-za\). The projectors of distinct
pairs commute. Their product \(\Pi_{\mathrm{vac}}\) is
finite on an input and projects to the common annihilator
kernel \(\operatorname{Vac}(M)\).

Taylor expansion gives
\[
 \begin{gathered}
 m=\sum_{\alpha,I}\frac{z^\alpha z_I}{\alpha!}
         \Pi_{\mathrm{vac}}(a_Ia^\alpha m),\\
 P_{\ell,\mathrm{gr}}\otimes_{S_\ell}\operatorname{Vac}(M)
       \xrightarrow{\sim}M.
 \end{gathered}
\]
Odd annihilators are ordered opposite to the creation monomial.
Applying projected derivatives proves uniqueness. Projecting
a lift of a vacuum vector proves exactness of
\(\operatorname{Vac}\). The differential preserves its kernel
and becomes \(D_P\otimes1+1\otimes d_{\mathrm{vac}}\).
This gives the same derived coefficient category on the
jointly supported modules. Fixing oscillator translation
makes the comparison translation-compatible. An arbitrary
extra covariant translation is instead an extra degree-zero
chain operator on the vacuum complex and must be retained.

This positive theorem uses full quartet-linear maps.
Forgetting to the boundary category is not full:
\([\Omega,\theta_0]=\eta_1\ne0\), and the scalar end has
no degree-two class mapping to the boundary residue.
Nor does primitive diagonal tensoring give a normalized
quartet tensor:
\[
 [\xi\otimes1+1\otimes\xi,\ x\otimes1+1\otimes x]
       =2\,1.
\]
The construction therefore does not identify a native
relative sewing product.

\section{Endomorphisms under coefficient and collision maps}
\label{sec:vol2-marked-base-change}

The endomorphisms of a marked object need not commute with
restriction of its coefficients. The comparison is induced
by evaluation and exists before it is known to be invertible.
Let \(R\to S\) be a map of commutative dg rings, \(B\) a
dg \(R\)-algebra, and \(P,Q\) left \(B\)-modules. Write
\(B_S=S\otimes_R^{\mathbb L}B\) and similarly for modules.
Extending evaluation and taking its adjoint gives
\begin{equation}\label{eq:vol2-marked-base-change}
 S\otimes_R^{\mathbb L}\operatorname{RHom}_B(P,Q)
 \longrightarrow
 \operatorname{RHom}_{B_S}(P_S,Q_S).
\end{equation}

\begin{proposition}\label{prop:vol2-perfect-marked-base-change}
If \(P\) is perfect over \(B\), the comparison
\eqref{eq:vol2-marked-base-change} is an equivalence for
every \(Q\). For \(Q=P\), it is an equivalence of unital
\(E_1\)-algebras. For two perfect modules over
\(R\)-algebras \(B_1,B_2\), the external-product comparison
\[
 \operatorname{REnd}_{B_1}(P_1)
       \otimes_R^{\mathbb L}\operatorname{REnd}_{B_2}(P_2)
 \longrightarrow
 \operatorname{REnd}_{B_1\otimes_R^{\mathbb L}B_2}
                (P_1\otimes_R^{\mathbb L}P_2)
\]
is also an equivalence. Both assertions retain a finite
diagram of marking arrows when all its source objects
are perfect.
\end{proposition}
\begin{proof}
For \(P=B\), \eqref{eq:vol2-marked-base-change} is the
identity on \(S\otimes_R^{\mathbb L}Q\). Both sides, as
functors of \(P\), take finite extensions to finite limits
and preserve retracts. Scalar extension preserves finite
limits in a stable category. Induction on finite cells
proves the assertion for every perfect \(P\).
For the external product, replace the target copies of
\(P_1,P_2\) by arbitrary \(Q_1,Q_2\). The comparison is the
identity for the two free source modules; apply the same
finite-cell argument in each source variable.
The maps are adjoint to evaluation, so they preserve
composition, units and successive base changes. Cross Hom
has the same proof, and the arrow end model uses a finite
homotopy fibre. This proves the marked assertion.
\end{proof}

Perfectness cannot simply be omitted. Let \(R=k[t]\),
\(S=R[t^{-1}]\) and \(P=\bigoplus_{n\ge0}Re_n\).
The endomorphism \(e_n\mapsto t^{-n}e_n\) of \(P_S\)
has no common denominator, so it is absent from
\(S\otimes_R\operatorname{End}_R(P)\).
There is a separate tensor defect. On
\(V=\bigoplus_{i\ge0}ke_i\), the endomorphism
\[
 T(e_i\otimes e_j)=\delta_{ij}e_i\otimes e_j
\]
does not belong to \(\operatorname{End}(V)\otimes
\operatorname{End}(V)\). A sum of \(r\) elementary tensors
would make every finite diagonal coefficient matrix have
rank at most \(r\), contradicting its \((r+1)\)-square
identity corner.

\subsection{The collision comparisons}

Let \(X\) be a smooth curve and \(C\) a connected
coaugmented factorizing coalgebra on its ordered collision
carrier. Write \(\mathscr R_I\) for the carrier with a finite
nonempty set of roots \(I\). Suppose
a marked bicomodule category \(\mathcal M_I\) and its
regular object \(\mathbb V_I\) have been specified, with
a closed action of the normalized coefficient category.
On a smooth chart \(Y\) of dimension \(d\), and for a
collision map \(f:Y'\to Y\), use
\[
 L_Y(M)=\omega_Y\otimes_{\mathcal O_Y}M[d],\qquad
 f^\dagger=L_{Y'}^{-1}f^!L_Y.
\]
This fixes the left-to-right normalization; in particular
\(f^\dagger\mathcal O_Y=\mathcal O_{Y'}\).
Define the internal end
\[
 E_I=\underline{\operatorname{RHom}}_{\mathcal M_I}
                   (\mathbb V_I,\mathbb V_I).
\]
Its enrichment, rather than a global scalar Hom, determines
the coefficient object.

If collision restriction takes the regular object to the
regular object and is linear over normalized pullback,
evaluation induces
\begin{equation}\label{eq:vol2-root-collision-end-map}
 c_\alpha:\delta_\alpha^\dagger E_I\longrightarrow E_{I'}
 \qquad(\alpha:I\twoheadrightarrow I').
\end{equation}
For a partition \(\mathcal P=(I_1|\cdots|I_r)\), write
\(j_{\mathcal P}:U_{\mathcal P}\to\mathscr R_I\) for the
separated open. Let \(E_{I,\mathcal P}^{\mathrm{loc}}\)
be the end of the regular object in the restricted marked
category. Restriction and external-product evaluation give
the span
\begin{equation}\label{eq:vol2-separated-end-span}
 j_{\mathcal P}^\dagger E_I
 \xrightarrow{\ b_{\mathcal P}\ }E_{I,\mathcal P}^{\mathrm{loc}}
 \xleftarrow{\ e_{\mathcal P}\ }
 \left.\boxtimes_a E_{I_a}\right|_{U_{\mathcal P}}.
\end{equation}
Neither arrow is inverted by its construction.

\begin{proposition}\label{prop:vol2-perfect-collision-criterion}
Suppose the entire marked category on each root chart is
presented as \(\operatorname{Mod}_{B_I}\), carrying
\(\mathbb V_I\) to a perfect \(P_I\). Suppose each collision
and separated-open restriction is the corresponding derived
extension of scalars, and the separated categories and regular
objects are the external tensor products of their blocks.
If these identifications preserve coefficient actions,
chart transitions, promotions of labels and every ordered
collision diagram, then \(c_\alpha,b_{\mathcal P},e_{\mathcal P}\)
are coherent equivalences preserving composition and evaluation.
\end{proposition}
\begin{proof}
Apply Proposition~\ref{prop:vol2-perfect-marked-base-change}
on each chart. The maps obtained there are the maps adjoint
to the same evaluations as
\eqref{eq:vol2-root-collision-end-map} and
\eqref{eq:vol2-separated-end-span}. The supplied coherent
presentations therefore identify them on overlaps and on
composite restrictions. Equivalences of underlying
coefficient complexes detect equivalences of the normalized
differential modules. They descend with their algebra and
evaluation structures.
\end{proof}

Perfectness at each fixed number of ordinary labels is
weaker than the premise. It neither identifies the restriction
functor of the entire marked category with scalar extension
nor commutes an infinite homotopy end with pullback.

Let \(\mathcal C^f\) be the category of factorizing coalgebras
on this carrier, with transverse tensor \(\otimes^!\)
and local unit \(J\). A transverse monoid is a unital
associative algebra object for this tensor.
The space \(\operatorname{Act}(F;C)\) consists of its
unital associative actions on \(C\), with all their
comparison homotopies.

\begin{theorem}[Conditional factorizing recognition]
\label{thm:vol2-factorizing-recognition}
Suppose the comparisons
\eqref{eq:vol2-root-collision-end-map} and
\eqref{eq:vol2-separated-end-span} are equivalences with all
mixed collision and factorization coherences. Suppose their
evaluations descend, preserve the local unit, and commute
with promotion of ordinary labels to roots. Then
\[
 Z_I=L_{\mathscr R_I}(E_I),\qquad Z_\varnothing=k,
\]
defines a factorizing transverse monoid with an action on \(C\).
It represents the full transverse action functor if, in
addition, evaluation induces equivalences
\begin{equation}\label{eq:vol2-native-action-recognition}
 \operatorname{Map}_{\operatorname{Alg}(\mathcal C^f)}(F,Z)
       \xrightarrow{\sim}\operatorname{Act}(F;C)
\end{equation}
for every actor \(F\) in the specified category
\(\mathcal C^f\).
\end{theorem}
\begin{proof}
The maps \(b_{\mathcal P}^{-1}e_{\mathcal P}\) give the
separated-product equivalences, and the \(c_\alpha\) give
collision equivalences. Their stipulated coherences are the
descent diagrams. Composition and the unit descend because
the comparisons preserve them; the evaluation hypothesis
supplies the action. Composing a map \(F\to Z\) with this
action constructs \eqref{eq:vol2-native-action-recognition}
on all simplices, naturally in \(F\). Its additional
equivalence assertion is exactly representability of the
full action functor.
\end{proof}

The additional assertion is substantive. The full pointed
and ordinary multiplicative action categories considered in
Chapter~\ref{chap:vol2-recipients} have explicit
nonfinality obstructions. A proposed restriction must explain
which actions or maps it excludes and why that restriction is
geometrically appropriate. The marked associative classifier
of Theorem~\ref{thm:vol2-fixed-mark-classifier} remains valid
in its arrow category. Identifying that category with the
supported factorizing action category is a separate map.
The two-object question of a pair of actions with no common
recipient also retains its own quantifier; an unequalizable
parallel pair does not by itself supply that pair of objects.

\section{Complete filtrations and formal comparison}
\label{sec:vol2-complete-comparison}

A comparison of finite coefficients becomes a comparison of
formal objects only when the limit and its operations are
controlled. Let \(C\) be a complex with descending
subcomplex filtration
\(C=F^0C\supset F^1C\supset\cdots\).
Call it complete and separated here when the natural map
\[
 C\longrightarrow\varprojlim_n C/F^nC
\]
is an isomorphism degree by degree. The differential is
continuous for this filtration. Set
\(\operatorname{gr}^nC=F^nC/F^{n+1}C\).

\begin{lemma}\label{lem:vol2-complete-filtered-comparison}
Let \(f:C\to D\) preserve complete separated filtrations
of this form. If every associated-graded map
\(\operatorname{gr}^nf\) is a quasi-isomorphism, then \(f\)
is a quasi-isomorphism. The conclusion is an equivalence
of structured algebras if \(f\) was already a continuous
map of those structured algebras.
\end{lemma}
\begin{proof}
Give the cone of \(f\) the componentwise filtration. Its
associated graded is acyclic. For a cone cocycle \(z\),
acyclicity of \(\operatorname{gr}^0\) gives \(b_0\) with
\(z-db_0\in F^1\). This error is closed. Inductively,
acyclicity of \(\operatorname{gr}^n\) gives
\(b_n\in F^n\) such that
\[
 z-d(b_0+\cdots+b_n)\in F^{n+1}.
\]
Completeness defines \(b=\sum_{n\ge0}b_n\), and continuity
with separatedness gives \(db=z\). Thus the cone is
acyclic. No new operation is supplied by this argument;
it proves equivalence for the continuous structured map
assumed in the statement.
\end{proof}

The existence of a compatible map is a different problem.
For example, let \(\mathfrak g\) be a complete dg Lie
algebra with
\([F^i\mathfrak g,F^j\mathfrak g]\subset F^{i+j}\mathfrak g\).
A deformation is a Maurer--Cartan element
\(\alpha\in F^1\mathfrak g^1\):
\[
 d\alpha+\tfrac12[\alpha,\alpha]=0.
\]
If this equation holds modulo \(F^n\), \(n\ge1\), its error
\(r_n\in F^n\mathfrak g^2\) is closed modulo \(F^{n+1}\).
Indeed the Bianchi identity
\(dr_n+[\alpha,r_n]=0\) makes the second term one
filtration step higher. A correction \(u_n\in F^n\mathfrak g^1\)
changes the error by \(du_n\) modulo \(F^{n+1}\).
Thus the next obstruction lies in
\(H^2(\operatorname{gr}^n\mathfrak g)\). When it vanishes,
choices differ by degree-one classes; their gauge
identifications involve degree-zero classes.
Completeness sums actual compatible corrections. It does
not force these obstruction or automorphism groups to vanish.

This distinction preserves the Hamiltonian comparisons of
the preceding chapters. A complete cochain map can be a
quasi-isomorphism while its proposed brace refinement remains
unconstructed. A selected higher coefficient can change under
transfer. A full Massey coset, a nonzero action-cone class or
a residual gauge stabilizer requires its own invariant
calculation.

\section{Holomorphic descent with a specified comparison}
\label{sec:vol2-marked-holomorphic-descent}

The Hamiltonian Lie-cochain model and the symplectic
polyvector model have different carriers. To formulate the
second, let \((S,\omega)\) be a holomorphic symplectic
surface and \(R=\mathbb C\llbracket\hbarq\rrbracket\).
Write \(\pi=\omega^{-1}\). A chosen formal Poisson tensor
\(\pi_\hbar=\hbarq\pi+O(\hbarq^2)\), with
\([\pi_\hbar,\pi_\hbar]=0\), defines
\[
 \mathcal P_S=
 \bigl(\operatorname{PV}_S\llbracket\hbarq\rrbracket,
                  [\pi_\hbar,-]\bigr).
\]
Polyvector degree is its cohomological degree; the Schouten
bracket has degree \(-1\), so this differential has degree one.
Here \(\operatorname{PV}_S^p=\bigwedge^pT_S\).
The choice \(\pi_\hbar=\hbarq\pi\) is a particular formal
class, not the only quantization datum.

A Weyl algebroid \(\mathcal W\) is specified by local
formal Weyl algebras, equivalences on overlaps and their
coherent cocycle data. A Fedosov presentation uses formal
fibre coordinates and a flat resolution differential, locally
written
\[
 D_F=\nabla-\delta+\hbarq^{-1}[\gamma,-].
\]
Here \(\nabla\) is the chosen local symplectic connection,
\(\gamma\) is a Weyl-valued one-form, and \(\delta\) is the
fibre Koszul differential. The
commutator is divisible by \(\hbarq\), so the displayed
operator need not invert the coefficient parameter.
Flatness is an equation on \(\gamma\) and the prescribed
curvature data. In the holomorphic setting these resolution
and descent data must be supplied, for example through
compatible Dolbeault models; they are not obtained by
assuming a global holomorphic symplectic connection exists.

Dolgushev's smooth formality construction uses Fedosov
resolutions and gives an \(L_\infty\) comparison of
polyvectors with polydifferential operators
\cite{DolgushevCov}*{Theorems 2--4}. That result alone
does not give a holomorphic homotopy-\(E_2\) comparison
with every marking and cyclic structure. We state precisely
the data used below.

Fix a cover, all its multiple intersections, and continuous
homotopy-\(E_2\) resolutions of \(\mathcal P_S\) and
\(\mathcal C_{\mathcal W}\), the relative Hochschild
cochain sheaf of \(\mathcal W\). Suppose a homotopy coherent
map of their descent diagrams
\begin{equation}\label{eq:vol2-fedosov-comparison-data}
 \Phi_\bullet:\mathcal P_\bullet\longrightarrow
                            \mathcal C_{\mathcal W,\bullet}
\end{equation}
has been constructed. Its leading maps are required to be
quasi-isomorphisms for complete separated positive
filtrations as in
Lemma~\ref{lem:vol2-complete-filtered-comparison}.
The coefficient completion, chosen formality refinement
and formal Poisson class are fixed throughout.

\begin{theorem}[Descent of a supplied formality comparison]
\label{thm:vol2-fedosov-marked-descent}
Under these hypotheses, \(\Phi_\bullet\) descends to an
equivalence of the complete homotopy-\(E_2\) sheaves and
induces
\[
 \mathbf R\Gamma(S,\mathcal P_S)
       \simeq \mathbf R\Gamma(S,\mathcal C_{\mathcal W}).
\]
If the chosen continuous \(R\)-linear Morita theory
identifies the latter with the centre of
\(\operatorname{Perf}(\mathcal W)\), this is the corresponding
global centre comparison. Coherent markings, cyclic maps
and state maps descend whenever
\(\Phi_\bullet\) is supplied in their diagram category.
\end{theorem}
\begin{proof}
The filtered lemma makes the map an equivalence on every
member of the descent diagram. A homotopy limit sends an
objectwise equivalence of diagrams to an equivalence.
Taking the descent totalization therefore gives the sheaf
equivalence, and derived global sections give the displayed
map. Limits of algebra objects carry their operations
objectwise, so the homotopy-\(E_2\) structure is preserved.
The Morita assertion uses its stated centre model.
The same limit argument in a diagram category retains
markings and all their comparison homotopies.
\end{proof}

A change of Fedosov data gives the same conclusion when
a compatible gauge equivalence of the chosen resolutions
and its descent homotopies are supplied. An arbitrary
change of the Fedosov one-form is not such a hypothesis:
it can change the formal curvature class.
Likewise, equivariance requires a coherent equivariant
choice of \eqref{eq:vol2-fedosov-comparison-data}.
These qualifications apply on a K3 surface as well.
Its holomorphic symplectic form supplies \(\pi\), not all
of the displayed comparison data.

There is also a formal algebraic realization in two
topological directions. For each holomorphic open, the
specified \(E_2\)-algebra has its factorization realization
on framed two-disks. The framed structure may be promoted
to an oriented one when the required rotation-equivariant
structure is supplied. If these assignments satisfy the
required cosheaf descent in the holomorphic variable, they
define a mixed factorization theory on
\(\mathbb R^2_{\mathrm{top}}\times S\).
Theorem~\ref{thm:vol2-fedosov-marked-descent} then realizes
the supplied centre equivalence in those topological
directions. Neither this realization nor Morita invariance
alone identifies a universal physical boundary category.

For the Hamiltonian BF action
\[
 S_{\mathrm{BF}}(A,B)=
 \int_{\mathbb R^2\times S}
 \left\langle B,(d+\bar\partial)A+\tfrac12[A,A]\right\rangle,
\]
the first comparison is a map from its Hamiltonian Chevalley
observables to the chosen symplectic polyvector complex.
It must preserve the operations and coefficient completion.
Gauge fixing and a coupled quantum master equation are
additional requirements for a quantum factorization theory.
When \(S\) is a complex surface, this spacetime has six real
dimensions. The three-dimensional half-space theory of
Chapter~\ref{ch:vol2-native-bulk-boundary} has a different
geometric input.

\section{Sewing with a prepared state}
\label{sec:vol2-marked-state-sewing}

A reconstruction can preserve an algebra and lose its state.
Let \(B\) be a dg \(R\)-algebra, \(M\) a \(B\)-module,
\(\nu:R\to M\) a chosen closed preparation in \(R\)-complexes,
and \(s:M\to R\) a closed \(R\)-linear state map with
\(s\nu=1\). An equivalence of such data consists of the
algebra and module comparisons, a homotopy carrying the
preparation, a homotopy carrying \(s\), and compatibility
of these homotopies with \(s\nu=1\).
Neither the unit of \(B\) nor an augmentation of \(B\)
specifies \(s\).

To express sewing, fix boundary algebras \(B_0,B_1,B_2\),
a \(B_0\)--\(B_1\) bimodule \(M_{01}\), and a
\(B_1\)--\(B_2\) bimodule \(M_{12}\). The sewn object
\(M_{02}\) is a \(B_0\)--\(B_2\) bimodule.
A sewing map has domain
\[
 \mu_{012}:
 M_{01}\widehat\otimes^{\mathbb L}_{B_1}M_{12}
 \longrightarrow M_{02},
\]
where the completed balanced tensor has been independently
defined on the allowed objects. Preparations must obey
\(\mu_{012}(\nu_{01}\otimes\nu_{12})\simeq\nu_{02}\).
If \(\sigma_{012}\) is the supplied contraction functional
on that balanced tensor, state compatibility is the
specified homotopy
\begin{equation}\label{eq:vol2-marked-sewing-state}
 s_{02}\mu_{012}\simeq\sigma_{012}.
\end{equation}
For three regions these homotopies must agree with the
associativity comparison; higher regions require their
higher coherence. A tensor of scalar functionals descends
to a balanced tensor only when it satisfies the actual
balancing relation. This must be checked in the state model.

\begin{proposition}\label{prop:vol2-sewing-transport}
Suppose two such systems are related by equivalences of
their full diagrams of boundary algebras, bimodules,
markings, state maps and sewing maps, including every
coherence. Then all finite sewn states agree under those
equivalences. If the systems are diagrams of derived
complete objects and their comparisons commute with the
coefficient quotients, the same is true after the prescribed
derived completion.
\end{proposition}
\begin{proof}
For two regions the assertion is
\eqref{eq:vol2-marked-sewing-state} transported through
the given square. Inductively compose this square with
the next sewing map. The associativity homotopies identify
the different parenthesizations, and their higher
coherences identify the resulting paths of identifications.
Thus the full finite diagram, including its final state
map, is carried by the equivalence. At the complete level,
apply the derived inverse limit to the compatible
coefficient diagrams. It preserves their equivalences
and homotopies.
\end{proof}

This proposition requires states on the completed domain in
question. Finite Gaussian integration constructs the
normalization and balanced contractions in
Chapter~\ref{ch:vol2-gaussian-elimination}; the determinant
cannot be removed from that state comparison.
The infinite-mode and stabilized states have the specific
analytic domains of
Chapter~\ref{ch:vol2-infinite-mode-states}.
A formal coefficient equivalence by itself supplies
neither positivity nor a Laplace prescription, and
coefficientwise existence does not establish continuity
on a larger state completion.

\begin{problem}[Local mapping and physical reconstruction]
\label{prob:vol2-marked-physical-reconstruction}
Construct a stable boundary category with homotopy-invariant
cross-object local mapping objects, independently defined
relative sewing and evaluation against both boundary
algebras. The construction must retain the coefficient
lines, local units, preparations and regular, localized,
supported and vacuum collision terms. Determine its
comparison with the marked associative action classifier
and with independently defined bulk observables.

For the Ramond junction, this requires deciding how the
local residue \(q^2\Omega\) and its cone occur in that
category, including the derivative contact
\eqref{eq:vol2-quartet-derivative-contact}.
A category containing the acyclic localization \(Q_\ell\)
cannot have unchanged polynomial local actors on every
diagonal and still be derived invariant. Within the
specified even-boundary derived category, no full stable
sector containing the junction and its shifts can require
every object to lift to the original quartet.
An admissible alternative must construct its mapping and
sewing operations and prove its comparison, rather than
infer them from the end algebra of one object.
\end{problem}
